\chapter{Requisitos Funcionais e Reengenharia}
\label{sec:reqfunc}
\index{Requisitos Funcionais}

\section{Levantamento de Requisitos Funcionais}
\label{sec:levantamento}

\subsection{Inventário das Funcionalidades Existentes}
\label{subsec:inventário}

\subsubsection{Funções do Sistema}
\label{ssubsec:funcoes}

\begin{quadro}[ht]
\begin{center}
\small
\captionsetup{justification=centering, labelsep=endash, width=0.75\linewidth}
\caption{Inventário das funções do sistema}
\label{quadro:funcoes}
\begin{threeparttable}
\begin{tabularx}{\textwidth}{|l|l|X|} \hline
\textbf{Categoria} & \textbf{Funcionalidade} & \textbf{Descrição Técnica} \\ \hline
Cálculos & RLM Simples (Baseline) & Regressão Linear Múltipla padrão sem técnicas de ensemble \\ \hline
Cálculos & RLM + Bagging & Ensemble com amostragem \textit{com} reposição (bootstrap) \\ \hline
Cálculos & RLM + Subagging + CV & Ensemble com amostragem \textit{sem} reposição + otimização por validação cruzada 5-Fold \\ \hline
Configuração & Parâmetro $m$ (Bags) & Define número de amostras/ensembles para o método Bagging \\ \hline
Configuração & Parâmetro $k$ (Subagging) & Define fração ou quantidade para amostragem no Subagging \\ \hline
Visualização & Métricas de Desempenho & Exibição de R², RMSE, MAE calculados no conjunto de teste \\ \hline
Visualização & Gráficos Comparativos & \textit{Scatter plots} (previsto $\times$ real), resíduos e comparação entre modelos \\ \hline
Organização & Abas de Navegação & Separação lógica: ``Arquivos'', ``Configuração'', ``Resultados'', ``Gráficos'' \\ \hline
Exportação & \textit{Não identificada} & \textbf{Obs.:} sistema não apresenta botões de download de resultados \\ \hline
\end{tabularx}
\begin{tablenotes}[flushleft]
\footnotesize
\item Fonte:~\oautor
\end{tablenotes}
\end{threeparttable}
\end{center}
\end{quadro}

\newpage
\subsubsection{Entradas Necessárias}
\label{ssubsec:entradas}

\begin{quadro}[ht]
\begin{center}
\small
\captionsetup{justification=centering, labelsep=endash, width=0.75\linewidth}
\caption{Entradas necessárias para execução do sistema}
\label{quadro:entradas}
\begin{threeparttable}
\begin{tabularx}{\textwidth}{|l|X|X|} \hline
\textbf{Tipo} & \textbf{Nome/Parâmetro} & \textbf{Descrição} \\ \hline
Arquivo & \texttt{Xcal.txt} & Matriz de espectros de calibração (variáveis independentes) \\ \hline
Arquivo & \texttt{Ycal.txt} & Vetor de valores de referência para calibração \\ \hline
Arquivo & \texttt{Xteste.txt} & Matriz de espectros para predição/teste \\ \hline
Arquivo & \texttt{Yteste.txt} & Vetor de valores reais para validação do conjunto de teste \\ \hline
Parâmetro & $m$ (nº de bags) & Inteiro $\geq 1$. Usado apenas nos métodos de ensemble \\ \hline
Parâmetro & $k$ (Subagging) & Decimal $(0, 1]$. Usado apenas no método Subagging \\ \hline
Config. & Validação Cruzada & Fixo: 5-Fold. Parâmetro interno de otimização do modelo \\ \hline
\end{tabularx}
\begin{tablenotes}[flushleft]
\footnotesize
\item Fonte:~\oautor
\end{tablenotes}
\end{threeparttable}
\end{center}
\end{quadro}

\subsubsection{Saídas Esperadas}
\label{ssubsec:saidas}

\begin{quadro}[ht]
\begin{center}
\small
\captionsetup{justification=centering, labelsep=endash, width=0.75\linewidth}
\caption{Saídas esperadas do sistema}
\label{quadro:saidas}
\begin{threeparttable}
\begin{tabularx}{\textwidth}{|l|l|X|} \hline
\textbf{Tipo} & \textbf{Conteúdo} & \textbf{Formato de Exibição} \\ \hline
Numérico & R², RMSE, MAE (teste) & Tabela ou \textit{cards} destacados \\ \hline
Gráfico & Previsto vs. Real & \textit{Scatter plot} com reta de ajuste \\ \hline
Gráfico & Resíduos da Predição & Gráfico de dispersão ou histograma \\ \hline
Comparativo & Desempenho entre modelos & Tabela comparativa ou gráfico de barras \\ \hline
Feedback & Mensagens de erro/validação & Alertas em destaque na interface \\ \hline
Status & Indicador de processamento & \textit{Spinner} ou barra de progresso \\ \hline
\end{tabularx}
\begin{tablenotes}[flushleft]
\footnotesize
\item Fonte:~\oautor
\end{tablenotes}
\end{threeparttable}
\end{center}
\end{quadro}

\newpage

\subsection{Classificação por Prioridade}
\label{subsec:prioridade}

\subsubsection{Essenciais (\textit{Must Have})}
\label{ssubsec:essenciais}

\begin{quadro}[ht]
\begin{center}
\small
\captionsetup{justification=centering, labelsep=endash, width=0.75\linewidth}
\caption{Requisitos essenciais do sistema}
\label{quadro:essenciais}
\begin{threeparttable}
\begin{tabularx}{\textwidth}{|l|X|X|} \hline
\textbf{ID} & \textbf{Requisito} & \textbf{Justificativa} \\ \hline
E01 & Upload e validação dos 4 arquivos & Dados são pré-requisito para qualquer cálculo estatístico \\ \hline
E02 & Execução da RLM Simples (baseline) & Referência obrigatória para comparação de desempenho \\ \hline
E03 & Cálculo de métricas básicas (R², RMSE, MAE) & Mínimo necessário para avaliação quantitativa da qualidade \\ \hline
E04 & Validação dimensional (X vs Y) & Previne erros de execução e resultados inválidos \\ \hline
E05 & Exibição clara dos resultados & Permite decisão técnica fundamentada em evidências \\ \hline
\end{tabularx}
\begin{tablenotes}[flushleft]
\footnotesize
\item Fonte:~\oautor
\end{tablenotes}
\end{threeparttable}
\end{center}
\end{quadro}

\subsubsection{Desejáveis (\textit{Should Have})}
\label{ssubsec:desejaveis}

\begin{quadro}[ht]
\begin{center}
\small
\captionsetup{justification=centering, labelsep=endash, width=0.75\linewidth}
\caption{Requisitos desejáveis do sistema}
\label{quadro:desejaveis}
\begin{threeparttable}
\begin{tabularx}{\textwidth}{|l|X|X|} \hline
\textbf{ID} & \textbf{Requisito} & \textbf{Benefício} \\ \hline
D01 & Gráficos interativos (zoom, \textit{tooltip}) & Facilita exploração visual e apresentação de resultados \\ \hline
D02 & Exportação de resultados (CSV, PDF, PNG) & Permite documentação e reprodutibilidade da análise \\ \hline
D03 & \textit{Tooltips} explicativos para $m$, $k$ & Reduz curva de aprendizado para usuários não especialistas \\ \hline
D04 & Pré-processamento básico (normalização) & Melhora robustez com dados espectrais reais \\ \hline
D05 & Histórico de execuções na sessão & Permite comparar configurações sem reenviar arquivos \\ \hline
\end{tabularx}
\begin{tablenotes}[flushleft]
\footnotesize
\item Fonte:~\oautor
\end{tablenotes}
\end{threeparttable}
\end{center}
\end{quadro}

\newpage

\subsubsection{Extras (\textit{Could Have})}
\label{ssubsec:extras}

\begin{quadro}[ht]
\begin{center}
\small
\captionsetup{justification=centering, labelsep=endash, width=0.75\linewidth}
\caption{Requisitos extras para fases futuras}
\label{quadro:extras}
\begin{threeparttable}
\begin{tabularx}{\textwidth}{|l|X|X|} \hline
\textbf{ID} & \textbf{Requisito} & \textbf{Cenário de Uso} \\ \hline
X01 & Autenticação e projetos privados & Ambientes corporativos ou pesquisa com dados sensíveis \\ \hline
X02 & Integração com APIs externas & Pipelines automatizadas de quimiometria \\ \hline
X03 & Suporte a formatos (.csv, .xlsx, .mat) & Flexibilidade para ingestão de diferentes fontes \\ \hline
X04 & Otimização automática de hiperparâmetros & Reduz necessidade de conhecimento técnico avançado \\ \hline
X05 & Dashboard comparativo multi-rodada & Análise estatística de estabilidade dos ensembles \\ \hline
\end{tabularx}
\begin{tablenotes}[flushleft]
\footnotesize
\item Fonte:~\oautor
\end{tablenotes}
\end{threeparttable}
\end{center}
\end{quadro}

\newpage
\subsection{Fluxos de Uso}
\label{subsec:fluxos}

\subsubsection{Fluxo Principal}
\label{ssubsec:fluxo-principal}

%----------------------------------------------------------------
\begin{figure}[htbp]
\centering
\captionsetup{justification=centering, labelsep=endash, width=0.75\linewidth}
\caption{Fluxo principal de execução da análise espectral}
\label{fig:fluxo-principal}
\begin{tikzpicture}[
  node distance=0.6cm and 1.5cm,
  startstop/.style={rectangle, rounded corners, minimum width=3.2cm, minimum height=0.8cm,
                    align=center, draw=black, line width=0.8pt, fill=gray!15, font=\small\sffamily\bfseries},
  process/.style={rectangle, minimum width=5.8cm, minimum height=0.85cm,
                  align=center, draw=black, line width=0.8pt, fill=blue!5, font=\small\sffamily},
  decision/.style={diamond, aspect=2.2, minimum width=2.8cm, minimum height=1.0cm,
                   align=center, draw=black, line width=0.8pt, fill=orange!15, font=\small\sffamily\bfseries, inner sep=1.5pt},
  arrow/.style={thick,->,>=Stealth}
]

% Nós principais
\node[startstop] (start) {Início};
\node[process, below=of start] (upload) {Upload: Xcal, Ycal, Xteste, Yteste};
\node[decision, below=of upload] (valid) {Arquivos\\válidos?};
\node[process, below=of valid] (config) {Habilita aba ``Configuração''};
\node[process, below=of config] (params) {Define $m$ e $k$ (se ensemble)};
\node[process, below=of params] (exec) {Executa Análise com Otimização CV};
\node[process, below=of exec] (progress) {Exibe indicador de processamento};
\node[process, below=of progress] (calc) {Processa: RLM $\rightarrow$ Bagging $\rightarrow$ Subagging+CV};
\node[process, below=of calc] (results) {Exibe aba ``Resultados'' com métricas};
\node[process, below=of results] (charts) {Navega para ``Gráficos''};

% CORREÇÃO: Alterado de 'startstop' para 'process' com preenchimento cinza para ter espaço para o texto
\node[process, below=of charts, fill=gray!15, font=\small\sffamily\bfseries] (end) {Fim: Interpretação técnica};

% Nó de tratamento de erro
\node[process, right=of valid, fill=red!10, minimum width=4cm] (error) {Alerta: ``Verifique\\formato/dimensões''};

% Setas
\draw[arrow] (start) -- (upload);
\draw[arrow] (upload) -- (valid);
\draw[arrow] (valid) -- node[right, font=\small\sffamily] {Sim} (config);
\draw[arrow] (config) -- (params);
\draw[arrow] (params) -- (exec);
\draw[arrow] (exec) -- (progress);
\draw[arrow] (progress) -- (calc);
\draw[arrow] (calc) -- (results);
\draw[arrow] (results) -- (charts);
\draw[arrow] (charts) -- (end);

\draw[arrow] (valid) -- node[above, font=\small\sffamily] {Não} (error);

% CORREÇÃO: Saída pelo topo do bloco de erro (error.north)
\draw[arrow] (error.north) |- (upload);

\end{tikzpicture}
\caption*{\centering Fonte: \oautor}
	\label{fig:figura14}
\end{figure}
%---------------------------------------------------------------

\newpage
%--------------------------------------------------
\subsubsection{Fluxos Alternativos e Pontos Críticos}
\label{ssubsec:pontos-criticos}

\begin{quadro}[ht]
\begin{center}
\small
\captionsetup{justification=centering, labelsep=endash, width=0.75\linewidth}
\caption{Tratamento de exceções e usabilidade}
\label{quadro:excecoes}
\begin{threeparttable}
\begin{tabularx}{\textwidth}{|l|X|X|} \hline
\textbf{Cenário} & \textbf{Comportamento do Sistema} & \textbf{Mitigação Proposta} \\ \hline
Arquivo inválido & Alerta vermelho + destaque no campo & Validação dimensional \textit{fail-fast} antes do processamento \\ \hline
Parâmetros fora do intervalo & Validação em tempo real + mensagem & \textit{Sliders} com valores recomendados e limites explícitos \\ \hline
\textit{Timeout} de processamento & Mensagem orientativa de otimização & Barra de progresso com etapas descritivas \\ \hline
Navegação sem executar & Persistência de uploads na sessão & Estado mantido em \texttt{st.session\_state} \\ \hline
Resultados não exportáveis & Impede documentação externa & Botão ``Baixar Relatório'' com métricas + gráficos \\ \hline
\end{tabularx}
\begin{tablenotes}[flushleft]
\footnotesize
\item Fonte:~\oautor
\end{tablenotes}
\end{threeparttable}
\end{center}
\end{quadro}

\subsection{Requisitos Não Funcionais}
\label{subsec:nao-funcionais}

\begin{quadro}[ht]
\begin{center}
\small
\captionsetup{justification=centering, labelsep=endash, width=0.75\linewidth}
\caption{Especificação de requisitos não funcionais}
\label{quadro:rnf}
\begin{threeparttable}
\begin{tabularx}{\textwidth}{|l|X|X|} \hline
\textbf{Categoria} & \textbf{Especificação} & \textbf{Critério de Aceite} \\ \hline
Performance & Tempo de cálculo (RLM/Ensemble) & RLM $<$ 10\,s; Ensemble $<$ 2\,min (CV 5-Fold) \\ \hline
Performance & Renderização de gráficos & $<$ 3\,s após conclusão do cálculo \\ \hline
Usabilidade & Acessibilidade (WCAG 2.1 AA)(1) & Contraste, \textit{labels}, navegação por teclado \\ \hline
Usabilidade & Clareza textual e recuperação de erro & Mensagens orientam correção, não apenas apontam falha \\ \hline
Segurança & Isolamento de sessão e sanitização & Dados não vazam; validação de extensão/conteúdo \\ \hline
Escalabilidade & Crescimento de usuários/dados & Container com \textit{load balancer}; processamento em \textit{batch} \\ \hline
\end{tabularx}
\begin{tablenotes}[flushleft]
\footnotesize \item Fonte:~\oautor
\item (1) Padrão W3C de acessibilidade web. Garante que conteúdo seja perceptível, operável, compreensível e robusto para pessoas com deficiência. Nível AA = conformidade legal recomendada (contraste 4.5:1, teclado, alt text, legendas, ARIA).
\end{tablenotes}
\end{threeparttable}
\end{center}
\end{quadro}

\newpage
\section{Reengenharia pelo Modelo Clássico}
\label{sec:reengenharia}

\subsection{Diagnóstico e Projeto Lógico}
\label{subsec:diag-logico}

\begin{quadro}[ht]
\begin{center}
\small
\captionsetup{justification=centering, labelsep=endash, width=0.75\linewidth}
\caption{Diagnóstico do código legado e arquitetura proposta}
\label{quadro:diag}
\begin{threeparttable}
\begin{tabularx}{\textwidth}{|l|X|X|} \hline
\textbf{Aspecto} & \textbf{Estado Atual (\texttt{app.py})} & \textbf{Solução na Reengenharia} \\ \hline
Arquitetura & Monolítica (UI + lógica + \textit{parsing}) & Camadas separadas: Apresentação, Negócio, Dados \\ \hline
Validação & \texttt{try/except} genérico & Função \texttt{validate\_inputs()} com regras explícitas \\ \hline
Performance & Laços \texttt{for} em CV, sem cache & \texttt{@st.cache\_data}, vetorização \texttt{numpy}, \textit{seed} exposto \\ \hline
Tipagem/Doc & Ausente em funções críticas & \textit{Type hints} + \textit{docstrings} padrão NumPy \\ \hline
Reprodutibilidade & \texttt{seed=42} \textit{hardcoded} & Parâmetro configurável + registro em logs \\ \hline
\end{tabularx}
\begin{tablenotes}[flushleft]
\footnotesize
\item Fonte:~\oautor
\end{tablenotes}
\end{threeparttable}
\end{center}
\end{quadro}

\subsection{Estrutura Física e Implementação}
\label{subsec:fisico-impl}

\begin{quadro}[ht]
\begin{center}
\small
\captionsetup{justification=centering, labelsep=endash, width=0.75\linewidth}
\caption{Estrutura modular e responsabilidades}
\label{quadro:estrutura}
\begin{threeparttable}
\begin{tabularx}{\textwidth}{|l|X|} \hline
\textbf{Módulo/Arquivo} & \textbf{Função Específica} \\ \hline
\texttt{app.py} & Camada de apresentação (apenas UI + roteamento de abas) \\ \hline
\texttt{src/data\_loader.py} & \textit{Parsing} de \texttt{.txt}, validação dimensional, \textit{dataclasses} \\ \hline
\texttt{src/models.py} & Implementação de RLM, Bagging e Subagging (vetorizado) \\ \hline
\texttt{src/metrics.py} & Cálculo de RMSE, MAE e R² com tipos explícitos \\ \hline
\texttt{src/optimizers.py} & Validação cruzada 5-Fold e busca de hiperparâmetros \\ \hline
\texttt{tests/} & Suíte de testes unitários (\texttt{pytest}) para cada módulo \\ \hline
\end{tabularx}
\begin{tablenotes}[flushleft]
\footnotesize
\item Fonte:~\oautor
\end{tablenotes}
\end{threeparttable}
\end{center}
\end{quadro}

\newpage
\subsection{Testes, Validação e Rastreabilidade}
\label{subsec:testes-rtm}

\begin{quadro}[ht]
\begin{center}
\small
\captionsetup{justification=centering, labelsep=endash, width=0.75\linewidth}
\caption{Estratégia de testes e critérios de aceite}
\label{quadro:testes}
\begin{threeparttable}
\begin{tabularx}{\textwidth}{|l|X|X|} \hline
\textbf{Tipo de Teste} & \textbf{Escopo} & \textbf{Meta de Cobertura} \\ \hline
Unitário & Funções isoladas (\texttt{parse}, \texttt{metrics}) & $\geq$ 80\% (\texttt{pytest-cov}) \\ \hline
Integração & Fluxo completo: Upload $\rightarrow$ CV $\rightarrow$ Gráficos & 100\% dos caminhos críticos \\ \hline
Performance & Tempo de resposta para \textit{datasets} típicos & RLM $<$ 10\,s; Ensemble $<$ 2\,min \\ \hline
Aceite (UAT)(1) & Validação por pesquisadores/alunos & SUS $\geq$ 75 (Bom/Excelente)(2) \\ \hline
\end{tabularx}
\begin{tablenotes}[flushleft]
%\footnotesize
\item \footnotesize{Fonte:~\oautor}
\item (1) Teste final executado pelo usuário/cliente para validar se o sistema atende aos requisitos de negócio em cenários reais. Aprovação formal (\textit{sign-off}) libera a solução para produção.
\item (2) Pontuação na System Usability Scale (questionário de 10 itens) que indica usabilidade percebida acima da média (percentil ~73+). Meta comum em projetos UX para validar que o produto é intuitivo, eficiente e satisfatório para os usuários.
\end{tablenotes}
\end{threeparttable}
\end{center}
\end{quadro}

\begin{quadro}[ht]
\begin{center}
\small
\captionsetup{justification=centering, labelsep=endash, width=0.75\linewidth}
\caption{Matriz de rastreabilidade de requisitos (RTM)}
\label{quadro:rtm}
\begin{threeparttable}
\begin{tabularx}{\textwidth}{|l|X|X|} \hline
\textbf{ID Requisito} & \textbf{Fase Clássica} & \textbf{Artefato de Validação} \\ \hline
RF01.1--01.3 & 1, 3, 5 & \texttt{data\_loader.py}, testes de upload \\ \hline
RF02.1--02.3 & 1, 2, 4 & \texttt{optimizers.py}, logs de execução \\ \hline
RF03.1--03.3 & 2, 3, 4 & \texttt{models.py}, testes unitários de métricas \\ \hline
RF04.1--04.5 & 2, 4, 5 & \texttt{metrics.py}, UAT visual (\textit{Plotly}) \\ \hline
RF05.1--05.3 & 1, 3, 5 & UI/UX, \textit{feedback} de erros, \texttt{st.progress} \\ \hline
\end{tabularx}
\begin{tablenotes}[flushleft]
\footnotesize
\item Fonte:~\oautor
\end{tablenotes}
\end{threeparttable}
\end{center}
\end{quadro}